\section{Belief-State Controller 作为研究主线的理论依据与问题定义}

\subsection{从当前系统到控制问题的提出}

前文已经说明，当前系统并不缺少工作流执行机制。无论是显式有限状态机、阶段化蛋白设计流程，还是人工介入、恢复链和事件审计，这些能力都表明系统已经具备较完整的多智能体运行时结构。换言之，当前系统面临的主要矛盾，已经不再是“能不能把流程跑起来”，而是“在一个高代价、长链路、可失败、可恢复的科学工作流中，系统究竟依据什么原则决定下一步如何执行”。

这个问题之所以重要，是因为蛋白设计这样的科学工作流并不是简单的问答任务或单步工具调用任务。其执行过程通常包含多个强耦合阶段，不同步骤的成本和失败代价显著不对称，而且错误往往会沿工作流向后传播。一个系统即便具备了候选生成、结构预测、质量门控和人工介入能力，如果它仍然无法回答“现在是否应该继续计算、切换路径、请求人工介入或提前终止”，那么它本质上仍然只是一个更复杂的执行平台，而不是一个具备独立算法主线的研究系统。

也正是在这一意义上，controller 才成为比“继续做生命科学助手”更合适的研究主线。它使系统的核心问题从“完成一个生命科学任务”转变为“如何在风险、成本、恢复空间和人工预算约束下控制科学工作流的推进方式”。这不仅更接近控制、调度和序贯决策问题，也更有助于切断与“生命科学助手”父集之间的包含关系，因为 controller 所研究的是更底层的执行治理机制，而不是具体的任务入口或交互界面。

\subsection{为什么选择 Belief-State Controller}

在 controller 的不同建模方式中，之所以特别考虑 \texttt{Belief-State Controller}，是因为当前系统最关键、也最难被直接观测的因素，恰恰都具有“隐状态”的特征。系统无法在执行之前明确知道某个任务到底有多难、当前候选是否已经足够可靠、当前局部失败究竟是偶发噪声还是结构性困难、剩余恢复空间是否足以支持继续自动执行，以及当前请求人工介入是否真正具有边际价值。换言之，系统面临的并不是一个完全可观测的决策环境，而是一个必须根据运行时证据不断更新判断的不完全信息环境。

如果采用静态评分或规则门控，这些问题往往会被简化为一组人工权重和阈值，例如“如果风险分高于某一阈值则请求人工介入，否则继续执行”。这种做法虽然简单，但会把“难度”退化为固定标签，把“控制”退化为启发式分支选择。相反，\texttt{Belief-State Controller} 的价值在于，它允许系统将“难度、风险、证据充分性、恢复余量和人工介入价值”统一看作需要持续估计的隐状态，并根据新观测不断更新对这些隐状态的判断。

更重要的是，Belief State 在本课题中的含义需要明确说明。本文将“信念状态（belief state）”理解为：\textbf{在已有运行时观测和历史动作条件下，系统对当前执行环境中若干不可直接观测变量的联合估计。} 它不是单一的难度分数，也不是某个简单风险标签，而是一组关于当前工作流真实状态的结构化估计，例如当前任务难度、当前证据是否充分、当前恢复空间是否仍充足、以及当前人工介入是否具有价值等。也正因为如此，Belief State 并不是一个附加说明字段，而是控制器做出治理动作选择的状态基础。

\subsection{本文采用的问题定义}

在本课题中，\texttt{Belief-State Controller} 面对的对象不是单轮对话，也不是单次模型推理，而是一条高代价科学工作流。该工作流可以抽象为一个带有数据依赖与控制依赖的图结构，其中每个节点对应某一阶段的任务执行、质量判断、恢复操作或人工介入点。对于蛋白设计场景而言，这样的工作流至少包含候选生成、结构预测、质量门控、结构条件优化以及恢复等关键部分。

在任意时刻，系统都处在一个无法被完全直接观测的运行状态中。这个状态并不是单一的 \texttt{difficulty score}，而是一个结构化的隐状态集合，用于共同刻画当前任务的真实执行难度、当前证据是否充分、当前工具链是否稳定、当前恢复空间是否仍然充足、以及当前是否值得引入人工介入。系统所能获得的只是运行过程中的局部观测，例如候选结果的一致性和分歧程度、结构预测输出及其质量指标、质量门控的拒绝原因、最近一次 \texttt{patch} 或 \texttt{replan} 的结果、已经消耗的时间与资源，以及此前人工决策留下的轨迹。controller 的任务，正是依据这些异构运行时观测，持续更新系统对于当前隐状态的信念状态，并据此在若干治理动作之间做出选择。

由此，本文中的 \texttt{Belief-State Controller} 可以被理解为一个面向科学工作流的部分可观测治理控制器。它不直接生成最终任务结果，而是以工作流的运行方式为控制对象，在成本、风险、时间和人工预算约束下，决定工作流下一步应如何推进、恢复、暂停或终止。这个定义具有两层意义：一方面，它使“难度”和“风险”真正进入系统决策闭环；另一方面，它也使课题从一个面向生命科学的助手系统，转化为一个面向高代价科学工作流的控制问题。

\subsection{相关研究现状及其启示}

围绕这一问题，现有研究已经提供了若干重要启示，但尚未直接给出与本课题完全重合的解决方案。

第一类相关工作是不确定性感知规划。例如 UoT\parencite{uot} 提出在信息不足时让 language agent 根据不确定性主动进行信息搜集，这说明 uncertainty-aware planning 已经是被认可的研究方向，也证明“不确定性可以进入决策闭环”这一思路是成立的。与本课题相比，这类工作更多关注信息寻求与提问行为，而较少涉及高代价工作流中的阶段化恢复、人工介入和治理动作，因此它为本研究提供的是方法论启发，而不是现成答案。

第二类相关工作是成本感知路由与 \texttt{test-time} 资源控制。BEST-Route\parencite{best-route} 将任务难度与质量阈值纳入模型路由和采样预算分配，EcoAct\parencite{ecoact} 则从经济性角度讨论 agent 系统中何时注册何种动作或工具。它们共同说明，在大模型系统中，根据难度和成本动态分配计算资源已经成为清晰的研究问题。这与本课题高度相关，因为 \texttt{Belief-State Controller} 同样需要根据状态估计决定是否继续投入更多执行代价。但不同之处在于，这些工作通常停留在模型路由或工具使用层面，而本文所关注的是工作流层面的恢复、人工介入和执行治理。

第三类相关工作是长程价值引导与搜索控制。例如 QLASS\parencite{qlass} 通过逐步价值估计引导 agent 在复杂任务中做更好的中间选择，LATS\parencite{lats} 则通过树搜索统一推理、行动与规划。这些工作表明，大模型 agent 的关键不只在于最终输出，更在于中间决策质量以及对未来收益的估计能力。对本课题而言，其重要意义在于说明：在长链路任务中，控制“下一步如何走”本身就是研究重点。然而，这类工作更偏向搜索效率和推理路径优化，而不是面向高代价工作流的治理控制。

第四类相关工作是多智能体系统失败分析与结构化设计。Why Do Multi-Agent LLM Systems Fail?\parencite{whyfail} 指出，多智能体系统的失败很大一部分并非来自模型本身，而来自角色失配、接口不清、验证不足以及终止条件设计不良。Which Agent Causes Task Failures and When?\parencite{whichandwhen} 进一步说明，多智能体链式执行中的失败归因并不容易。与此同时，MetaAgent\parencite{metaagent} 这类工作也显示，显式状态机等结构已经开始被用于更系统化地组织多智能体系统。这些研究共同提供了重要背景：在多智能体系统中，“如何治理执行过程”已经成为一个独立问题，而不是简单附属于任务求解过程。

综合这些文献可以看出，现有研究分别从不确定性、成本约束、价值引导、失败分析和结构化设计等角度接近了本课题，但尚未把以下几个因素统一起来：高代价科学工作流、显式恢复动作、人工介入动作、结构化工作流状态以及运行时信念状态更新。也正因为如此，\texttt{Belief-State Controller} 仍然具有明确的研究空间。

\subsection{本节小结}

将当前选题进一步收束为 \texttt{Belief-State Controller}，并不是因为这一概念“听起来更高级”，而是因为它同时满足三个条件：第一，它最能回应当前系统的真正缺口，即运行时控制依据的缺失；第二，它最能把难度、风险、工作量、恢复与人工介入统一为一个明确的算法问题；第三，它在现有研究中有清晰的相关脉络可依托，又没有被已有工作完全覆盖。因此，它不是一个凭空引入的新名词，而是当前条件下最有希望同时兼顾算法性、独立性、系统兼容性与时间可行性的研究方向。
